\section{유한체의 크기와 존재}
\label{sec:finite-field-construction}

\begin{learninggoal}
유한체의 원소 수가 소수의 거듭제곱이어야 함을 증명한다. 다항식 $X^q-X$의 근들을 이용하여 원소가 $q=p^n$개인 체를 구성하고, 그 체의 유일성을 설명한다.
\end{learninggoal}

\begin{definition}[유한체]
원소의 개수가 유한한 체를 \emph{유한체}(finite field)라 한다. 원소가 $q$개인 유한체가 존재할 때 이를 보통 $\F_q$로 쓴다.
\end{definition}

유한체의 원소 수는 임의의 자연수가 될 수 없다. 표수와 벡터공간 차수가 가능한 크기를 완전히 제한한다.

\begin{theorem}[유한체의 가능한 크기]
$F$가 유한체이면 어떤 소수 $p$와 정수 $n\ge1$이 존재하여
\[
  |F|=p^n
\]
이다. 여기서 $p=\charac F$이고 $n=[F:\F_p]$이다.
\end{theorem}

\begin{proof}
유한체의 소체도 유한하므로 $F$의 표수는 어떤 소수 $p$이고 소체는 $\F_p$이다. $F$는 유한집합이므로 $\F_p$ 위의 벡터공간 차수
\[
  n=[F:\F_p]
\]
도 유한하다. 기저를 하나 택하면 각 원소는 $n$개의 $\F_p$-좌표로 유일하게 표현된다. 각 좌표에는 $p$개의 선택이 있으므로
\[
  |F|=p^n.
\]
\end{proof}

\begin{pitfall}
원소가 $p^n$개인 체 $\F_{p^n}$와 잉여류환 $\Z/p^n\Z$는 $n>1$일 때 전혀 다른 대상이다. $\Z/p^n\Z$에는 영인수가 있으므로 체가 아니지만, $\F_{p^n}$는 모든 영이 아닌 원소가 역원을 갖는다.
\end{pitfall}

\subsection{$X^q-X$가 유한체를 기억한다}

이제 $q=p^n$이라 하자. 유한체의 모든 원소는 하나의 다항식의 근으로 동시에 포착된다.

\begin{proposition}
$F$가 원소가 $q=p^n$개인 유한체이면 모든 $a\in F$에 대해
\[
  a^q=a
\]
이다. 따라서 $F$는 $\F_p$ 위 다항식
\[
  X^q-X
\]
의 분해체이다.
\end{proposition}

\begin{proof}
$a=0$이면 자명하다. $a\ne0$이면 유한군 $F^\times$의 위수는 $q-1$이므로 라그랑주 정리에 의해
\[
  a^{q-1}=1.
\]
따라서 $a^q=a$이다. 즉 $F$의 $q$개 원소가 모두 $X^q-X$의 근이다. 이 다항식의 차수도 $q$이므로 이들이 모든 근이며, $F$는 모든 근으로 생성되는 분해체이다.
\end{proof}

\begin{theorem}[유한체의 존재와 유일성]
소수 $p$와 정수 $n\ge1$을 택하고 $q=p^n$이라 하자. 그러면 원소가 $q$개인 체가 존재한다. 또한 원소가 $q$개인 두 체는 $\F_p$를 고정하는 동형에 대해 서로 동형이다.
\end{theorem}

\begin{proofidea}
대수적 폐포 $\overline{\F_p}$ 안에서 $X^q-X$의 모든 근을 모은다. 이 근들의 집합이 사칙연산에 닫혀 있음을 보이면 원하는 체를 얻는다.
\end{proofidea}

\begin{proof}
$\overline{\F_p}$에서
\[
  f=X^q-X
\]
를 생각하자. 표수 $p$에서 $q=0$이므로
\[
  f'=qX^{q-1}-1=-1.
\]
따라서 $f$는 중근을 갖지 않고 정확히 $q$개의 서로 다른 근을 갖는다. 그 근들의 집합을
\[
  S=\{a\in\overline{\F_p}:a^q=a\}
\]
라 하자.

$a,b\in S$이면 $q=p^n$이므로 반복된 프로베니우스 공식에 의해
\[
  (a\pm b)^q=a^q\pm b^q=a\pm b,
  \qquad
  (ab)^q=a^qb^q=ab.
\]
또한 $b\ne0$이면
\[
  (b^{-1})^q=(b^q)^{-1}=b^{-1}.
\]
따라서 $S$는 체이고 원소 수는 $q$이다. 이를 $\F_q$라 쓴다.

유일성을 보자. 원소가 $q$개인 임의의 체는 앞의 명제에 의해 $X^q-X$의 $\F_p$ 위 분해체이다. 분해체는 기준체를 고정하는 동형까지 유일하므로 원소가 $q$개인 모든 체는 서로 동형이다.
\end{proof}

\begin{keyidea}
유한체 $\F_q$는 단순히 ``원소가 $q$개인 어떤 체''가 아니다. 고정된 대수적 폐포 안에서는
\[
  \F_q=\{a\in\overline{\F_p}:a^q=a\}
\]
로 볼 수 있다. 즉 $\F_q$는 프로베니우스의 고정점 전체이다.
\end{keyidea}

\subsection{작은 유한체의 구체적 구성}

\begin{example}[네 원소 체]
$X^2+X+1$은 $\F_2$에 근이 없으므로 기약이다. $\alpha^2+\alpha+1=0$인 $\alpha$를 첨가하면
\[
  \F_4=\F_2(\alpha)
  =\{0,1,\alpha,\alpha+1\}.
\]
관계식 $\alpha^2=\alpha+1$을 사용하면
\[
  \alpha^3=1.
\]
따라서 영이 아닌 세 원소는 모두 $X^3-1$의 근이고, 네 원소 모두 $X^4-X$의 근이다.
\end{example}

\begin{example}[여덟 원소 체]
$X^3+X+1$은 $\F_2$에서 기약이다. 한 근을 $\alpha$라 하면
\[
  \F_8=\F_2(\alpha),
  \qquad
  \alpha^3=\alpha+1.
\]
모든 원소는
\[
  a+b\alpha+c\alpha^2,
  \qquad a,b,c\in\F_2
\]
의 꼴로 유일하게 표현된다. 뒤에서 $\alpha$가 $\F_8^\times$ 전체를 생성함을 확인한다.
\end{example}

\begin{example}[아홉 원소 체]
$X^2+1$은 $\F_3$에 근이 없으므로 기약이다. $i^2=-1=2$인 원소를 첨가하면
\[
  \F_9=\F_3(i)
  =\{a+bi:a,b\in\F_3\}.
\]
이 체는 잉여류환 $\Z/9\Z$와 동형이 아니다. 예를 들어 $\Z/9\Z$에서는 $3\cdot3=0$이지만 $\F_9$에는 영이 아닌 영인수가 없다.
\end{example}

\begin{checkpoint}
\begin{enumerate}[label=(\alph*)]
  \item 원소가 $6$, $8$, $12$, $25$개인 체의 존재 가능성을 판정하라.
  \item $X^3+X+1$이 $\F_2[X]$에서 기약임을 확인하라.
  \item $\F_9$의 모든 원소가 $X^9-X$의 근인 이유를 설명하라.
\end{enumerate}
\end{checkpoint}
